% ==============================================================================
% config/elementos.tex
% Comandos de dados e leiaute dos elementos pré-textuais: capa, folha de rosto,
% folha de aprovação, resumos e epígrafe. Os dados em si ficam em dados.tex.
% ==============================================================================

% ------------------------------------------------------------------------------
% Instituição
% ------------------------------------------------------------------------------
\newcommand{\imprimiruniversidade}{Universidade Federal do Espírito Santo}
\newcommand{\imprimircentro}{Centro Universitário Norte do Espírito Santo}
\newcommand{\imprimirdepartamento}{Departamento de Computação e Eletrônica}

% ------------------------------------------------------------------------------
% Curso: \curso{cc} (Ciência da Computação) ou \curso{ec} (Engenharia de Computação)
% ------------------------------------------------------------------------------
% Em Engenharia de Computação, o TCC se chama Projeto de Graduação (nome da
% disciplina nas normas do curso). O comando \curso ajusta o curso, o grau e
% o nome do trabalho na folha de rosto e na folha de aprovação.
\newcommand{\imprimircurso}{Ciência da Computação}
\newcommand{\imprimirgrau}{Bacharel em Ciência da Computação}
\newcommand{\imprimirtipotcc}{Trabalho de Conclusão de Curso}
\newcommand{\curso}[1]{%
  \ifthenelse{\equal{#1}{ec}}{%
    \renewcommand{\imprimircurso}{Engenharia de Computação}%
    \renewcommand{\imprimirgrau}{Bacharel em Engenharia de Computação}%
    \renewcommand{\imprimirtipotcc}{Projeto de Graduação}%
  }{%
    \renewcommand{\imprimircurso}{Ciência da Computação}%
    \renewcommand{\imprimirgrau}{Bacharel em Ciência da Computação}%
    \renewcommand{\imprimirtipotcc}{Trabalho de Conclusão de Curso}%
  }%
}
\newcommand{\imprimircolegiado}{Colegiado do Curso de \imprimircurso}

% Natureza do trabalho (folha de rosto e folha de aprovação).
\newcommand{\imprimirnatureza}{%
  \imprimirtipotcc{} apresentado ao \imprimircolegiado{} do
  \imprimircentro{} da \imprimiruniversidade, como requisito parcial para
  obtenção do grau de \imprimirgrau.}
\preambulo{\imprimirnatureza}
\tipotrabalho{\imprimirtipotcc{} (Graduação)}

\local{São Mateus, ES}

% ------------------------------------------------------------------------------
% Dados complementares
% ------------------------------------------------------------------------------
% Subtítulo (opcional). Na capa e na folha de rosto aparece após dois-pontos.
\newcommand{\imprimirsubtitulo}{}
\newcommand{\subtitulo}[1]{\renewcommand{\imprimirsubtitulo}{#1}}
\newcommand{\imprimirtitulocompleto}{%
  \imprimirtitulo\ifdefempty{\imprimirsubtitulo}{}{: \imprimirsubtitulo}}

% Data de aprovação. Deixe vazio para gerar linhas em branco a preencher.
\newcommand{\imprimirdataaprovacao}{}
\newcommand{\dataaprovacao}[1]{\renewcommand{\imprimirdataaprovacao}{#1}}

% Banca examinadora: \membrobanca{Nome com titulação}{Instituição}{Função}
\makeatletter
\newcommand{\imprimirbanca}{}
\newcommand{\membrobanca}[3]{%
  \g@addto@macro\imprimirbanca{%
    \assinatura{\textbf{#1}\\#2\ifthenelse{\equal{#3}{}}{}{\\#3}}}%
}
\makeatother

% Palavras-chave (NBR 6028:2021): separadas por ponto e vírgula e
% finalizadas por ponto. Ex.: \palavraschave{otimização; heurística; mochila}
\newcommand{\palavraschave}[1]{%
  \par\vspace{\baselineskip}\noindent\textbf{Palavras-chave}: #1.}
\newcommand{\keywords}[1]{%
  \par\vspace{\baselineskip}\noindent\textbf{Keywords}: #1.}

% ------------------------------------------------------------------------------
% Capa (NBR 14724:2024, 4.1.1)
% ------------------------------------------------------------------------------
\renewcommand{\imprimircapa}{%
  \begin{capa}%
    \centering\SingleSpacing
    \includegraphics[width=3.3cm]{figuras/logo-ufes}\par
    \vspace{0.6cm}
    {\bfseries
      \MakeTextUppercase{\imprimircentro}\par
      \MakeTextUppercase{\imprimirdepartamento}\par
      \MakeTextUppercase{\imprimircolegiado}\par}
    \vfill
    {\large\imprimirautor\par}
    \vfill
    {\bfseries\Large\imprimirtitulocompleto\par}
    \vfill\vfill
    {\large\imprimirlocal\par
     \imprimirdata\par}
  \end{capa}%
}

% ------------------------------------------------------------------------------
% Folha de rosto (NBR 14724:2024, 4.2.1.1)
% ------------------------------------------------------------------------------
\renewcommand{\imprimirfolhaderosto}{%
  \begin{folhaderosto}%
    \centering\SingleSpacing
    {\large\imprimirautor\par}
    \vfill
    {\bfseries\Large\imprimirtitulocompleto\par}
    \vfill
    \noindent\hspace*{.45\textwidth}%
    \begin{minipage}{.55\textwidth}
      \justifying\imprimirnatureza

      \vspace{\baselineskip}
      \imprimirorientadorRotulo~\imprimirorientador
      \ifdefempty{\imprimircoorientador}{}{\par
        \imprimircoorientadorRotulo~\imprimircoorientador}
    \end{minipage}\par
    \vfill\vfill
    {\large\imprimirlocal\par
     \imprimirdata\par}
  \end{folhaderosto}%
}

% ------------------------------------------------------------------------------
% Folha de aprovação (NBR 14724:2024, 4.2.1.3)
% ------------------------------------------------------------------------------
\newcommand{\imprimirfolhadeaprovacao}{%
  \begin{folhadeaprovacao}%
    \centering
    {\large\imprimirautor\par}
    \vspace{1.2cm}
    {\bfseries\Large\imprimirtitulocompleto\par}
    \vspace{1.2cm}
    \noindent\hspace*{.45\textwidth}%
    \begin{minipage}{.55\textwidth}
      \justifying\imprimirnatureza
    \end{minipage}\par
    \vspace{1cm}
    \raggedright
    Aprovado em
    \ifdefempty{\imprimirdataaprovacao}%
      {\rule{1.2cm}{0.4pt} de \rule{3.5cm}{0.4pt} de \rule{1.6cm}{0.4pt}.}%
      {\imprimirdataaprovacao.}\par
    \vspace{0.8cm}
    \centering
    {\bfseries BANCA EXAMINADORA\par}
    \imprimirbanca
  \end{folhadeaprovacao}%
}

% \justifying para minipages (o memoir não define este comando).
\providecommand{\justifying}{%
  \rightskip=0pt plus 0pt\leftskip=0pt\parfillskip=0pt plus 1fil\relax}
